% !TEX root = ../main.tex
% !TEX program = lualatex
\documentclass[../main.tex]{subfiles}
\IfSubfilesClassLoaded{\externaldocument{\subfix{../build/main}}}{}
% =============================================================================
% File: 27I_DAWIA.tex
% =============================================================================
\begin{document}
\IfSubfilesClassLoaded{\chapter{EDO Study Guide}}{}
\section{Defense Acquisition Workforce Improvement Act (DAWIA)}

\subsection{Summary}
\begin{description}
  \item[\textbf{Purpose.}] \ac{dawia} professionalizes the \ac{dod} acquisition workforce through education, training, experience, and certification requirements~\autocite{edo-5-2-1-dawia-2025}.
  \item[\textbf{Positions and functional areas.}] Acquisition billets are positions where at least 50
  percent of duties are acquisition-related; the functional area and level in your orders or \ac{pd}
  drive certification requirements, and most EDO billets are coded for program management~\autocite{edo-5-2-1-dawia-2025}.
  \item[\textbf{Critical and key leadership roles.}] \acp{cap} and \acp{klp} carry tenure obligations and
  require \ac{apm} (or a waiver) to ensure experienced leadership continuity~\autocite{edo-5-2-1-dawia-2025}.
  \item[\textbf{Certification path.}] \ac{dawia} offers foundational, practitioner, and advanced levels;
  certification is requested in \ac{edacm} after meeting education, training, and experience standards~\autocite{edo-5-2-1-dawia-2025}.
  \item[\textbf{Continuous learning.}] Workforce members complete at least 80 hours of continuous learning
  every two years, with points requested through \ac{edacm} and no carryover between cycles~\autocite{edo-5-2-1-dawia-2025}.
\end{description}

\subsection{DAWIA Background and Back-to-Basics}
\begin{description}
  \item[\textbf{Origins.}] \ac{dawia} was enacted in 1990 (Public Law 101-510) to raise the professionalism of
  the acquisition workforce through standardized education, training, and experience requirements~\autocite{edo-5-2-1-dawia-2025}.
  \item[\textbf{Back-to-Basics.}] In 2020, \ac{usdas} launched the Back-to-Basics initiative to modernize
  \ac{dod} implementation, streamline certification, and emphasize lifelong learning~\autocite{edo-5-2-1-dawia-2025}.
\end{description}

\subsection{Acquisition Positions and Functional Areas}
\begin{description}
  \item[\textbf{Acquisition positions.}] Acquisition positions are the largest part of the \ac{awf} and require
  that at least half of assigned responsibilities are acquisition-related~\autocite{edo-5-2-1-dawia-2025}.
  \item[\textbf{Coding and level.}] For officers, acquisition coding appears in orders; for civilians, it is in the
  \ac{pd}. The functional area and level (often embedded in the \ac{aqd}) drive the required education,
  training, and experience standards~\autocite{edo-5-2-1-dawia-2025}.
  \item[\textbf{Functional areas.}] \ac{dawia} functional areas include program management, engineering and
  technical management, contracting, life-cycle logistics, business financial management and cost
  estimating, and test and evaluation~\autocite{edo-5-2-1-dawia-2025}.
\end{description}

\subsection{Critical Acquisition Positions and Key Leadership Positions}
\begin{description}
  \item[\textbf{Critical Acquisition Positions (CAPs).}] \acp{cap} are senior supervisory or managerial positions
  that require \ac{apm}, carry a minimum three-year tenure obligation, and for military billets typically
  require O-5 or above~\autocite{edo-5-2-1-dawia-2025}.
  \item[\textbf{Key Leadership Positions (KLPs).}] \acp{klp} are the most senior \ac{dawia} roles (less than one
  percent of the \ac{awf}), designated by \ac{asnrdanda} with \ac{usdas} attention, and usually filled by
  O-6 or GS-15 equivalents or higher with three-to-four year tenures~\autocite{edo-5-2-1-dawia-2025}.
\end{description}

\subsection{Certification Levels and Requirements}
\begin{description}
  \item[\textbf{Levels.}] \ac{dawia} certifications are offered at foundational, practitioner, and advanced levels
  across all functional areas~\autocite{edo-5-2-1-dawia-2025}.
  \item[\textbf{Orders-driven timelines.}] EDO acquisition orders typically require practitioner certification in
  the program management functional area within 60 months of reporting; advanced certification is
  expected for senior program management billets~\autocite{edo-5-2-1-dawia-2025}.
  \item[\textbf{Certification request.}] After completing education, training, and experience standards, submit a
  certification request in \ac{edacm}. Commands may certify O-6 and below when criteria are met, while
  \ac{datm} certifies Flag and \ac{ses} personnel~\autocite{edo-5-2-1-dawia-2025}.
\end{description}

\subsection{Training and Experience Credit}
\begin{description}
  \item[\textbf{Getting started.}] Use \ac{edacm} to set up a profile, search \ac{dau} training, and enroll in
  courses; new offerings typically open in August and early sign-up improves scheduling options~\autocite{edo-5-2-1-dawia-2025}.
  \item[\textbf{Funding.}] Selecting ``parent command will fund travel and per diem'' can improve course
  acceptance if the command agrees and \ac{datm} funding is available for priority students~\autocite{edo-5-2-1-dawia-2025}.
  \item[\textbf{Experience credit.}] Graduate education may be credited toward experience requirements when
  pursuing certification or \ac{apm}, with credit caps of 12 months (up to 18 months for subspecialty 5100P
  and 24 months for 5100N) and no reuse for multiple certifications~\autocite{edo-5-2-1-dawia-2025}.
\end{description}

\subsection{Acquisition Professional Membership}
\begin{description}
  \item[\textbf{Purpose.}] \ac{apm} provides a pool of qualified personnel for \acp{cap} and \acp{klp}; membership
  (or a waiver) is required for assignment to those positions~\autocite{edo-5-2-1-dawia-2025}.
  \item[\textbf{Eligibility.}] Requirements include a baccalaureate degree with at least 12 business-related
  credits, \ac{dawia} certification in any functional area, LCDR or above, and 48 months of documented
  acquisition experience (with limited graduate-study credit allowed)~\autocite{edo-5-2-1-dawia-2025}.
  \item[\textbf{Package.}] \ac{apm} requests are submitted in \ac{edacm} with a DAC application (NAVPERS-1301-86),
  \ac{edacm} transcript, college transcripts, and FITREPs when documenting non-acquisition coded time~\autocite{edo-5-2-1-dawia-2025}.
\end{description}

\subsection{Continuous Learning}
\begin{description}
  \item[\textbf{Requirement.}] Complete at least 80 hours of continuous learning every two years; points do not
  carry over to the next cycle, and members must request credit in \ac{edacm}~\autocite{edo-5-2-1-dawia-2025}.
\end{description}

\subsection{Quick Review}
\begin{itemize}
  \item \ac{dawia} sets education, training, experience, and certification standards for the \ac{dod} acquisition workforce~\autocite{edo-5-2-1-dawia-2025}.
  \item Acquisition billets are coded in orders or \ac{pd} and require certification tied to functional area and level~\autocite{edo-5-2-1-dawia-2025}.
  \item \acp{cap} and \acp{klp} require \ac{apm} and carry multi-year tenure obligations~\autocite{edo-5-2-1-dawia-2025}.
  \item Certification requests, training enrollment, and continuous learning credits flow through \ac{edacm}~\autocite{edo-5-2-1-dawia-2025}.
  \item \ac{apm} eligibility requires education, certification, rank, and documented acquisition experience~\autocite{edo-5-2-1-dawia-2025}.
\end{itemize}

\IfSubfilesClassLoaded{
  \RenewDocumentCommand{\entryname}{}{\textbf{\color{Modern} Acronym}}
  \RenewDocumentCommand{\descriptionname}{}{\textbf{\color{Modern} Definition}}
  \printnoidxglossary[
    type=\acronymtype,
    title=Acronyms,
    style=long-booktabs
  ]}{}
\end{document}
